% Chapter 1. Material: docs/requirements/audience.md, docs/research/literature-review.md,
% docs/research/preregistration.md (v0.4), reports/ (SRC-038). Topic and RQs — ADR-0052.
\chapter{Introduction}\label{ch:intro}

\section{Context and relevance}

Reading the Qur'an aloud is a daily practice for many Muslims, and correct pronunciation is part of the
practice itself. For Arabic speakers the main task is to learn the rules of recitation; for everyone else
the task starts much earlier, with the letters and the sounds behind them. In Russia and Central Asia this
path has a long and specific history. At the end of the nineteenth century the Jadid reformers replaced
the old method, in which pupils memorised the names of letters for years, with a sound-based method that
taught a letter as a sound and built syllables from it in small steps \cite{SRC-016,SRC-017}. Ahmad Hadi
Maksudi's primer \emph{al-Muallim al-Thani} (``The Second Teacher'') brought this method to the reading of the
Qur'an \cite{SRC-015,SRC-018}. More than a century later it is still printed by the Muslim board of Tatarstan and
officially distributed in Kazakhstan \cite{SRC-040,SRC-020}.

The primer has never been used without a teacher. A lesson follows the cycle of \emph{talaqqi}: the
teacher reads, the learner repeats, the teacher corrects every slip, and the learner moves on only after an
error-free reading \cite{SRC-021,SRC-058}. The cycle is established in the teaching tradition, but it ties every learner to a teacher's
time. Adults who study after work, live far from a mosque school or feel shy reading in front of others
practise alone between lessons, and nobody tells them when a \ar{ح} has turned into a Russian [x] or a long
vowel has become short.

Two developments make it reasonable to ask whether a machine could fill part of that gap. Automatic speech
recognition has reached the level of phonemes for Qur'anic Arabic: the open model Quran Muaalem recognises
phonemes together with ten articulatory attributes and reports a phoneme error rate of about half a percent
on expert recitations \cite{SRC-002}. And tutors that listen to reading aloud already work for children
learning to read their native language \cite{SRC-125,SRC-130,SRC-132}; automatic pronunciation feedback has a
medium positive effect, largest when it targets individual sounds \cite{SRC-025}.

The topic is relevant for research, for practice and for industry. For research, open phoneme-level models
for Qur'anic Arabic appeared only recently \cite{SRC-002}, and public benchmarks contain no learners of Arabic as
a second language; the benchmark authors themselves name collecting such data as future work
\cite{SRC-004,SRC-031}. For practice, the primer is still printed and officially distributed
\cite{SRC-040,SRC-020}, so a tool built around it serves an existing teaching tradition instead of asking it to
change. For the speech-technology and educational-software industry, the thesis addresses a problem every team
faces when it deploys speech recognition for a group the model was not trained on: how to adapt an open model on
small data collected inside a working product and how to evaluate it without leakage.

The thesis belongs to the field of study 09.04.01 ``Computer Science'', profile Software Engineering: its
artefacts are a dataset, a model and the software that collects and processes the data; requirements are written
as quality-attribute scenarios, the architecture is checked mechanically, and the training and evaluation
pipeline is built to be reproducible.

\section{Problem}

An off-the-shelf model does not fit a beginner for two reasons. The first is who is speaking. Muaalem was
trained only on male expert reciters. On recordings with real errors its phoneme error rate is ten times higher
than on its own test set \cite{SRC-002}. Public benchmarks for Qur'anic pronunciation contain fluent verse
recitation and no learners of Arabic as a second language \cite{SRC-004,SRC-031,SRC-043}. Recognisers produce
false alarms on accented speech \cite{SRC-026}. A Russian-speaking adult reading a single syllable slowly for
the first time is far from the model's training data, and nobody knows how often the model will be wrong about
such a learner or whether it can be adapted to them. The methods exist --- low-rank fine-tuning
\cite{SRC-111}, calibration and abstention \cite{SRC-112,SRC-117} --- but they need data from the new
population, and there are none.

The second reason is what the model sees. It scores one attempt, while a teacher cares about something else:
\emph{what the learner has already mastered}. Knowledge tracing can answer that question from a sequence of
attempts \cite{SRC-124}, including from noisy recogniser output \cite{SRC-125}, but recogniser errors measured
against teachers have not been built into such models, and nobody has built them for adults learning to read
Arabic script.

Both reasons come down to one thing: there is no dataset. Without it one cannot measure how wrong the model is
on beginners, cannot fix it, and cannot check whether it estimates mastery correctly. The supervisor puts it
this way: what can be measured can be solved.

\section{Aim and objectives}

The aim is \emph{to develop a dataset of primer-level reading by Russian-speaking beginners and an AI model that
detects reading errors in it and estimates from speech the mastery of each skill of} al-Muallim al-Thani,
\emph{and to measure how close the model comes to a teacher and what it reveals about how people learn.}

Objectives:
\begin{enumerate}[label=O\arabic*.]
  \item turn the primer into a digital curriculum in the book's own order, with every exercise traceable to a
    page and line of an edition, and derive a skill map from it;
  \item build a tutor as the data-collection tool that recognises phonemes, explains errors in Russian to a
    learner starting from zero and keeps a full decision trace for every attempt;
  \item collect the dataset: consent, including consent to model training, and pseudonyms; a closed pilot with
    20--40 learners over 4--6 weeks; phoneme-level annotation with error type by two teachers, blind for test
    learners; a teacher test at three points; a datasheet;
  \item adapt the recognition model on teacher corrections and the book's minimal pairs and evaluate it on a
    frozen test set of new learners;
  \item build a mastery model of the book's skills that accounts for the recogniser's measured errors, and
    validate it on next attempts and against the teacher test;
  \item describe learners' progress through the book's skills and answer the research questions with effect
    sizes and confidence intervals under a preregistered plan.
\end{enumerate}

\section{Research questions}\label{sec:rq}

\begin{rqbox}
\textbf{RQ-1 (data).} How reliably can teachers and the model distinguish reading errors from pronunciation
errors in beginners reading the book?\\[4pt]
\textbf{RQ-2 (how the AI learns).} How much does adaptation on teacher corrections and the book's minimal
pairs improve recognition for new learners, and how does this depend on the amount of data?\\[4pt]
\textbf{RQ-3 (measuring mastery).} Can recognition be used to estimate whether a learner has mastered each
skill of the book, and does adapting the AI improve this estimate?\\[4pt]
\textbf{RQ-4 (how people learn).} How do Russian-speaking adults progress through the book's skills: what
comes quickly, what slowly, what is forgotten?
\end{rqbox}

The questions form a chain. RQ-1 checks what can be relied on: whether the annotation is reliable and how far
the off-the-shelf model is from a teacher. RQ-2 measures how the model learns from people. RQ-3 links the AI's
learning to the person's: if the adapted model hears more accurately, it should also estimate mastery more
accurately. RQ-4 uses this estimate and the teacher test to describe how people learn. The idea of the thesis is
``one primer, two learners'': the book's lessons define a skill map along which both the AI and the person
learn. Twelve confirmatory hypotheses are listed in Section~\ref{sec:method-hyp}; everything else is reported as
exploratory.

\section{Approach}

The thesis follows design science research \cite{SRC-100,SRC-101}. It builds three artefacts, and each claim is
evaluated with a suitable method. The dataset with its collection tool (Artefact A) is evaluated by measuring
annotation reliability. The recognition model (Artefact B) is evaluated by a controlled experiment on a frozen
test set of new learners. The mastery model (Artefact C) is evaluated by how well it predicts next attempts and
how well it agrees with a blind teacher test, and learners' progress is described longitudinally. Requirements
are written as quality-attribute scenarios \cite{SRC-047}, metrics are derived with the Goal--Question--Metric
method \cite{SRC-048}, and the analysis plan and data split will be preregistered and fingerprinted before research data
collection and before any training. The system is built by the author with an AI agent under mechanical checks
that protect both the code and the experiments \cite{SRC-067}.

\section{Contributions}

\begin{enumerate}
  \item[\textbf{C1.}] \textbf{Dataset.} The first corpus of reading by adult Russian-speaking beginners tied to
    the skills of \emph{al-Muallim al-Thani}, collected with consent, with phoneme-level and error-type
    annotation by two teachers, measured reliability, a blind gold test set and a datasheet; released together
    with the fine-tuned weights if ethics and consent allow.
  \item[\textbf{C2.}] \textbf{A model that learns from teachers.} Small-data adaptation of an open model to the
    new population --- a calibrated verifier whose false-rejection rate is controlled and checked on held-out learners, and LoRA fine-tuning
    with augmentation from the book's minimal pairs --- with a learning curve and the measured contribution of
    each data source.
  \item[\textbf{C3.}] \textbf{Mastery estimation from speech.} A mastery model that accounts for recogniser
    errors measured against teachers, validated on next attempts and against a blind teacher test.
  \item[\textbf{C4.}] \textbf{Knowledge about learners.} The first learning curves of Russian-speaking adults
    over the skills of a primer: what comes quickly, what slowly, what is forgotten.
\end{enumerate}

The engineering result is a reproducible ``collect $\to$ annotate $\to$ train $\to$ evaluate'' pipeline with
replaceable model, database and messenger.

\section{Object of research and novelty}\label{sec:novelty}
% Требование Положения о ГИА (п. 6.9): во введении — объект исследования и новизна.
% Формулировки — из вкладов C1--C4 и обзора (глава 2); [утвердить] с руководителем.

The \emph{object of research} is a dataset of primer-level reading by Russian-speaking adult beginners and an
AI model that detects reading errors in it and estimates skill mastery; the Telegram tutor is the instrument
that collects the data, not the object. The \emph{subject of research} is how adapting the model on teacher
corrections changes the agreement of its error detection and mastery estimates with teachers.

The novelty of the thesis, as far as the review in Chapter~\ref{ch:review} shows, is as follows:
\begin{enumerate}
  \item the dataset (C1) is the first collection of reading by adult Russian-speaking beginners tied to the
    skills of a primer, and its annotation separates reading errors from pronunciation errors;
  \item the mastery model (C3) takes into account the recogniser's errors measured against teachers;
    knowledge tracing on noisy recogniser output exists, but not with errors measured this way and not for
    adults learning to read Arabic script;
  \item the open Qur'anic phoneme model is adapted to a population it was not trained on (C2) using
    teacher corrections and the book's minimal pairs as data sources, with the contribution of each source
    measured on a learning curve;
  \item the learning curves (C4) are the first description of how Russian-speaking adults progress through
    the skills of \emph{al-Muallim al-Thani}.
\end{enumerate}

\section{Scope and limitations}

The curriculum covers the whole book (pp.~10--85), but evaluation focuses on the letter and vowel lessons
(pp.~10--24): learners start there, and they hold the most errors and skills. Tajweed rules such as
\emph{ikhfa} and \emph{idgham}, other first languages, choosing the next card from estimated mastery and a
controlled experiment on explanations are left for future work. The model is adapted, not trained from scratch:
the pilot yields orders of magnitude less data than training from scratch needs. RQ-4 describes how people
learn with the system, but without a control group it does not prove that the system is what taught them. The
tutor is meant to complement a teacher; the thesis does not claim that it can replace one --- and the book itself
says it cannot.

\section{Structure of the thesis}

Chapter~\ref{ch:review} reviews the primer and its method, the learners, research on pronunciation feedback,
automatic assessment of Qur'anic recitation, methods for adapting speech models, reading tutors and knowledge
tracing, and software engineering with AI agents. Chapter~\ref{ch:method} sets out the methodology:
requirements, metrics, four studies, preregistration and threats to validity. Chapter~\ref{ch:impl} describes
the implementation --- from the curriculum and skill map to the collection tool, the recognition and mastery
models and the development harness. Chapter~\ref{ch:eval} presents and discusses the evaluation, and
Chapter~\ref{ch:concl} concludes with answers and directions for the PhD.
